% Quiz originali pagine 1-10: 40 quesiti con spiegazioni rapide verificate.
% Sostituire il file quiz_pagine_001_010.tex esistente con questo file.
% Quesiti prima simulazione 7 e 16: risposta singola non valida; mantenuti con soluzione non verificata.

\par\noindent\hfill{\color{cyan!65!black}\scriptsize Pag. originale 1/521}\par\smallskip
\begin{quiz}
La derivata della funzione $f(x)=\sqrt[3]{e^{4x}\cos^{8}x}$ \`e:
\begin{enumerate}[label=\alph*),leftmargin=1.3cm,itemsep=3pt,topsep=5pt]
\item $\dfrac{1}{3\sqrt[3]{e^{8x}\cos^{16}x}}(32e^{4x}\sin^{7}x)$
\item $\dfrac{1}{3\sqrt[3]{e^{8x}\cos^{16}x}}(e^{4x}\cos^{8}x-e^{4x}\sin^{7}x)$
\item $\dfrac{e^{4x}}{3\sqrt[3]{e^{8x}\cos^{16}x}}(\cos^{8}x-\sin^{7}x)$
\item $\dfrac43\sqrt[3]{e^{4x}\cos^{5}x}(\cos x+2\sin x)$
\item $\dfrac43\sqrt[3]{e^{4x}\cos^{5}x}(\cos x-2\sin x)$
\end{enumerate}
\end{quiz}
% SOLUZIONE: E
% SEZIONE_AM1: 4.1
% SPIEGAZIONE: Metodo rapido: poni $u=e^{4x}\cos^8x$ e usa $(\sqrt[3]{u})\prime=u\prime/(3u^{2/3})$. Poiche $u\prime=4e^{4x}\cos^7x(\cos x-2\sin x)$, si ottiene $f\prime(x)=\frac43\sqrt[3]{e^{4x}\cos^5x}(\cos x-2\sin x)$. Risposta E.

\begin{quiz}
Sia $z=6-i$. Allora $w=z^2-\overline z$ vale
\begin{enumerate}[label=\alph*),leftmargin=1.3cm,itemsep=3pt,topsep=5pt]\item $w=29-13i$\item $w=29-11i$\item $w=37-13i$\item $w=37-11i$\item nessuna delle altre risposte\end{enumerate}
\end{quiz}
% SOLUZIONE: A
% SEZIONE_AM1: 1.5
% SPIEGAZIONE: Calcola $z^2=(6-i)^2=35-12i$. Attenzione al coniugato: $\overline z=6+i$. Quindi $w=z^2-\overline z=29-13i$. Risposta A.

\begin{quiz}
$\displaystyle\lim_{x\to+\infty}\frac{\sin(1/x)-2}{x^2}=$
\begin{enumerate}[label=\alph*),leftmargin=1.3cm,itemsep=3pt,topsep=5pt]\item non esiste\item $0$\item $+\infty$\item $1$\item nessuna delle altre risposte \`e corretta\end{enumerate}
\end{quiz}
% SOLUZIONE: B
% SEZIONE_AM1: 3.1
% SPIEGAZIONE: Il numeratore $\sin(1/x)-2$ e limitato (anzi tende a $-2$), mentre $x^2\to+\infty$. Quindi il rapporto tende a $0$. Risposta B.

\begin{quiz}
Quanti sono gli zeri della funzione $f(x)=x^9+4x+3$?
\begin{enumerate}[label=\alph*),leftmargin=1.3cm,itemsep=3pt,topsep=5pt]\item uno solo positivo\item uno solo negativo\item $2$\item $9$\item $9$, uno solo dei quali negativo\end{enumerate}
\end{quiz}
% SOLUZIONE: B
% SEZIONE_AM1: 4.3
% SPIEGAZIONE: Deriva: $f\prime(x)=9x^8+4>0$ per ogni $x$, quindi la funzione e strettamente crescente e puo avere al massimo uno zero. Poiche $f(-1)=-2<0$ e $f(0)=3>0$, per continuita lo zero esiste ed e negativo. Risposta B.

\begin{quiz}
Sapendo che $f(0)=4$, $f'(0)=3$ e posto $h(x)=\sqrt{f(x)}$, quanto vale $h'(0)$?
\begin{enumerate}[label=\alph*),leftmargin=1.3cm,itemsep=3pt,topsep=5pt]\item $2$\item $0$\item $3$\item $\frac38$\item $\frac34$\end{enumerate}
\end{quiz}
% SOLUZIONE: E
% SEZIONE_AM1: 4.1
% SPIEGAZIONE: Regola della catena: $h\prime(x)=f\prime(x)/(2\sqrt{f(x)})$. In $0$: $h\prime(0)=3/(2\sqrt4)=3/4$. Risposta E.

\begin{quiz}
$\displaystyle\lim_{n\to+\infty}\left(\cos\frac{2n^2}{n^3+1}\right)^{\arctan n}=$
\begin{enumerate}[label=\alph*),leftmargin=1.3cm,itemsep=3pt,topsep=5pt]\item $1$\item $\frac\pi2$\item $e$\item $0$\item $\infty$\end{enumerate}
\end{quiz}
% SOLUZIONE: A
% SEZIONE_AM1: 2.2
% SPIEGAZIONE: La base tende a $\cos0=1$ perche $2n^2/(n^3+1)\to0$; l esponente tende a $\pi/2$, che e finito. Quindi il limite e $1^{\pi/2}=1$. Risposta A.

\par\noindent\hfill{\color{cyan!65!black}\scriptsize Pag. originale 2/521}\par\smallskip
\begin{quiz}
L'asintoto obliquo della funzione $g(x)=\dfrac{x^2+2x-1}{3x-2}$ per $x\to+\infty$ \`e
\begin{enumerate}[label=\alph*),leftmargin=1.3cm,itemsep=3pt,topsep=5pt]\item $y=\frac13x-\frac19$\item $y=\frac32x+\frac94$\item $y=\frac13x+\frac29$\item $y=\frac32x-\frac{15}4$\item $y=\frac13x+\frac23$\end{enumerate}
\end{quiz}
% SOLUZIONE: ?
% SEZIONE_AM1: 4.4
% SPIEGAZIONE: ATTENZIONE: il quiz originale e incoerente. Dividendo $x^2+2x-1$ per $3x-2$ si ottiene $g(x)=x/3+8/9+(7/9)/(3x-2)$. L'asintoto e quindi $y=x/3+8/9$, non presente tra le cinque risposte. La tabella originale indica C, ma C e $y=x/3+2/9$: non selezionare una risposta come matematicamente corretta.

\begin{quiz}
$f(x)=\cos x$ \`e soluzione dell'equazione differenziale:
\begin{enumerate}[label=\alph*),leftmargin=1.3cm,itemsep=3pt,topsep=5pt]\item $(y')^2-y^2=1$\item $y'''+y=0$\item $y''-y=0$\item $(y')^2+y^2=1$\item $y''-2y=0$\end{enumerate}
\end{quiz}
% SOLUZIONE: D
% SEZIONE_AM1: 6.1
% SPIEGAZIONE: Per $y=\cos x$ vale $y\prime=-\sin x$. Per l identita trigonometrica $\sin^2x+\cos^2x=1$, si ha $(y\prime)^2+y^2=1$. Risposta D.

\begin{quiz}
$\displaystyle\lim_{x\to4/5}\frac{1-\cos(5x-4)}{(x-4/5)^2}=$
\begin{enumerate}[label=\alph*),leftmargin=1.3cm,itemsep=3pt,topsep=5pt]\item $\frac12$\item $1$\item $\frac52$\item nessuna delle altre risposte\item $\frac{25}2$\end{enumerate}
\end{quiz}
% SOLUZIONE: E
% SEZIONE_AM1: 3.1
% SPIEGAZIONE: Poni $t=5x-4=5(x-4/5)$. Usa $1-\cos t\sim t^2/2$ per $t\to0$: il rapporto e equivalente a $25(x-4/5)^2/[2(x-4/5)^2]=25/2$. Risposta E.

\begin{quiz}
$\displaystyle\lim_{n\to+\infty}\sqrt[n]{4^n+2^n}=$
\begin{enumerate}[label=\alph*),leftmargin=1.3cm,itemsep=3pt,topsep=5pt]\item $1$\item $+\infty$\item $e$\item $2$\item $4$\end{enumerate}
\end{quiz}
% SOLUZIONE: E
% SEZIONE_AM1: 2.2
% SPIEGAZIONE: Raccogli $4^n$ dentro la radice: $(4^n+2^n)^{1/n}=4(1+2^{-n})^{1/n}\to4$, poiche il fattore finale tende a $1$. Risposta E.

\begin{quiz}
Qual \`e la derivata prima della funzione $f(x)=\log(\cos x)$?
\begin{enumerate}[label=\alph*),leftmargin=1.3cm,itemsep=3pt,topsep=5pt]\item $f'(x)=\frac{\sin x}{\cos x}$\item $f'(x)=\tan x+c$\item $f'(x)=-\tan x$\item $\frac{-\sin x}{\log(\cos x)}$\item $-\sin x\cos x$\end{enumerate}
\end{quiz}
% SOLUZIONE: C
% SPIEGAZIONE: Derivata del logaritmo composto: $(\log u)\prime=u\prime/u$. Con $u=\cos x$ si ha $f\prime(x)=-\sin x/\cos x=-\tan x$ dove $\cos x>0$. Risposta C.

\begin{quiz}
La parte principale della funzione $f(x)=4x+\sin x+x^{6/5}$ per $x\to0$ \`e
\begin{enumerate}[label=\alph*),leftmargin=1.3cm,itemsep=3pt,topsep=5pt]\item $4x$\item $x^{6/5}$\item $x$\item $5x$\item $5x+x^{6/5}$\end{enumerate}
\end{quiz}
% SOLUZIONE: D
% SPIEGAZIONE: Per $x\to0$, $\sin x=x+o(x)$ e $x^{6/5}=o(x)$ per $x>0$; quindi $4x+\sin x+x^{6/5}=5x+o(x)$. La parte principale e $5x$. Risposta D. Nota: $x^{6/5}$ va inteso nel dominio reale appropriato.

\par\noindent\hfill{\color{cyan!65!black}\scriptsize Pag. originale 3/521}\par\smallskip
\begin{quiz}
Sia $f:[-2,2]\to\mathbb R$, $f(x)=\begin{cases}-x^5&x\ne0\\a&x=0.\end{cases}$
$f$ assume tutti i valori compresi tra il massimo e il minimo della funzione se:
\begin{enumerate}[label=\alph*),leftmargin=1.3cm,itemsep=3pt,topsep=5pt]\item $a=0$\item $\forall a\in[-2^5,2^5]$\item $a=2^5$\item $\forall a\in\mathbb R$\item nessuna delle altre risposte\end{enumerate}
\end{quiz}
% SOLUZIONE: A
% SPIEGAZIONE: Per $x\ne0$, $-x^5$ percorre tutti i valori di $[-32,32]$ tranne $0$. Per non lasciare un buco nell immagine occorre $f(0)=a=0$. Con $a=0$ la funzione e continua e assume tutti i valori intermedi. Risposta A.

\begin{quiz}
Sia $f$ una funzione continua e strettamente monotona, con $f(2)=7$, $f'(2)=\frac17$, $f'(7)=\frac12$. Allora $(f^{-1})'(7)=$
\begin{enumerate}[label=\alph*),leftmargin=1.3cm,itemsep=3pt,topsep=5pt]\item $2$\item $7$\item non \`e possibile saperlo\item $\frac17$\item $\frac12$\end{enumerate}
\end{quiz}
% SOLUZIONE: B
% SPIEGAZIONE: Derivata dell inversa: $(f^{-1})\prime(f(a))=1/f\prime(a)$, se $f\prime(a)\ne0$. Qui $f(2)=7$ e $f\prime(2)=1/7$, dunque $(f^{-1})\prime(7)=7$. Risposta B.

\begin{quiz}
$\displaystyle\lim_{x\to+\infty}x^2\left(\sqrt{1+\frac4x}-\cos\frac1x\right)=$
\begin{enumerate}[label=\alph*),leftmargin=1.3cm,itemsep=3pt,topsep=5pt]\item $+\infty$\item $0$\item $-\frac32$\item $1$\item $2$\end{enumerate}
\end{quiz}
% SOLUZIONE: A
% SPIEGAZIONE: Basta il termine dominante: $\sqrt{1+4/x}=1+2/x+O(1/x^2)$ e $\cos(1/x)=1-1/(2x^2)+O(1/x^4)$. La differenza e $2/x+O(1/x^2)$; moltiplicando per $x^2$ si ottiene $2x+O(1)\to+\infty$. Risposta A.

\begin{quiz}
L'insieme $A=\left\{-\frac{n}{n+5}:n=0,1,2,\ldots\right\}\cup\{-3\}$ ha
\begin{enumerate}[label=\alph*),leftmargin=1.3cm,itemsep=3pt,topsep=5pt]\item massimo, ma non minimo\item minimo, ma non massimo\item \`e illimitato\item $\inf(A)=-3$\item massimo e minimo\end{enumerate}
\end{quiz}
% SOLUZIONE: ?
% SPIEGAZIONE: Calcola gli estremi: $-n/(n+5)$ vale $0$ per $n=0$ e tende a $-1$ senza raggiungerlo. Aggiungendo $-3$ si ha $\\max A=0$, $\\min A=-3$ e $\\inf A=-3$. Quindi sono vere SIA D SIA E: domanda a risposta multipla erroneamente presentata come risposta singola.

\begin{quiz}
Sia data la funzione $f:\mathbb R\to\mathbb R$ tale che, per $x\to+\infty$, $f(x)\sim\log x$. Allora:
\begin{enumerate}[label=\alph*),leftmargin=1.3cm,itemsep=3pt,topsep=5pt]\item $e^{f(x)}\sim x$\item $\lim_{x\to+\infty}(f(x)-\log x)$ esiste\item $f(x)\,x\sim x^2$\item nessuna delle altre risposte \`e corretta\item $f(x)\log x=o(x^2)$ per $x\to+\infty$\end{enumerate}
\end{quiz}
% SOLUZIONE: E
% SPIEGAZIONE: Da $f(x)\sim\log x$ segue $f(x)/\log x\to1$, dunque $f(x)=O(\log x)$. Quindi $f(x)\log x=O((\log x)^2)=o(x^2)$. Risposta E. Le altre non seguono: per esempio $f(x)=\log x+\sin x$ rende $e^{f(x)}/x=e^{\sin x}$ senza limite.

\par\noindent\hfill{\color{cyan!65!black}\scriptsize Pag. originale 4/521}\par\smallskip
\begin{quiz}
La parte immaginaria del numero complesso $w=\dfrac{1+3i}{4-2i}$ \`e:
\begin{enumerate}[label=\alph*),leftmargin=1.3cm,itemsep=3pt,topsep=5pt]\item $3$\item $3i$\item $-\frac32$\item $\frac7{10}$\item $\frac7{10}i$\end{enumerate}
\end{quiz}
% SOLUZIONE: D
% SPIEGAZIONE: Moltiplica per il coniugato del denominatore: $w=(1+3i)(4+2i)/(4^2+2^2)=(-2+14i)/20=-1/10+(7/10)i$. La parte immaginaria e il numero reale $7/10$, non $7i/10$. Risposta D.

\begin{quiz}
Sia $f(x)=\sinh x$. Allora
\begin{enumerate}[label=\alph*),leftmargin=1.3cm,itemsep=3pt,topsep=5pt]\item $f^{(n)}(0)=0\iff n=4k$, $k\in\mathbb N$\item $f^{(n)}(0)=0\iff n=2k$, $k\in\mathbb Z$\item $f^{(n)}(0)=0\iff n=2k+1$, $k\in\mathbb N$\item $f^{(n)}(\pi)=0\iff n=2k$, $k\in\mathbb N$\item $f^{(n)}(0)=0\iff n=2k$, $k\in\mathbb N$\end{enumerate}
\end{quiz}
% SOLUZIONE: E
% SPIEGAZIONE: Le derivate di $\sinh x$ alternano $\cosh x$ e $\sinh x$. In $0$, $\sinh0=0$ e $\cosh0=1$: si annullano esattamente le derivate di ordine pari (incluso $0$). Risposta E.

\begin{quiz}
Sia data una funzione $f$ che soddisfa $-3^x\le f(x)\le3^x$, $f(2n+1)=3^{2n+1}$, $f(2n)=-3^{2n}$. Allora sicuramente
\begin{enumerate}[label=\alph*),leftmargin=1.3cm,itemsep=3pt,topsep=5pt]\item $\lim_{x\to+\infty}f(x)=\infty$\item $\lim_{x\to+\infty}f(x)=+\infty$\item $\lim_{x\to+\infty}f(x)=-\infty$\item $\lim_{x\to+\infty}f(x)=0$\item $\lim_{x\to+\infty}f(x)$ non esiste\end{enumerate}
\end{quiz}
% SOLUZIONE: E
% SPIEGAZIONE: Osserva le due sottosuccessioni: $f(2n)=-3^{2n}\to-\infty$ e $f(2n+1)=3^{2n+1}\to+\infty$. Poiche hanno comportamenti opposti, $\lim_{x\to+\infty}f(x)$ non esiste, neppure come limite infinito. Risposta E.

\par\noindent\hfill{\color{cyan!65!black}\scriptsize Pag. originale 6/521}\par\smallskip
\begin{quiz}
Sia $f:\mathbb R\to\mathbb R$ una funzione derivabile, tale che $f(1)=3$ e $f(4)=8$. Allora esiste $c\in(1,4)$ tale che
\begin{enumerate}[label=\alph*),leftmargin=1.3cm,itemsep=3pt,topsep=5pt]\item $f'(c)=0$\item $f'(c)=\frac35$\item $f'(c)=\frac53$\item $f(c)=\frac35$\item $f'(c)=1$\end{enumerate}
\end{quiz}
% SOLUZIONE: C
% SPIEGAZIONE: Teorema di Lagrange: una funzione derivabile su $\\mathbb R$ e continua su $[1,4]$. Esiste $c\\in(1,4)$ con $f'(c)=[f(4)-f(1)]/(4-1)=(8-3)/3=5/3$. Risposta 3 (C). Metodo rapido: differenza dei valori divisa per differenza degli argomenti.

\begin{quiz}
Il limite $\displaystyle\lim_{x\to+\infty}\frac{\cos x-x^2+e^{-2x}}{\sin x+e^{-3x}+x^2}$ vale
\begin{enumerate}[label=\alph*),leftmargin=1.3cm,itemsep=3pt,topsep=5pt]\item $1$\item $+\infty$\item $-\infty$\item $-1$\item $0$\end{enumerate}
\end{quiz}
% SOLUZIONE: D
% SPIEGAZIONE: Dividi numeratore e denominatore per $x^2$: $[\\cos x/x^2-1+e^{-2x}/x^2]/[\\sin x/x^2+e^{-3x}/x^2+1]$. I termini con seno, coseno ed esponenziali tendono a zero: resta $-1/1=-1$. Risposta 4 (D).

\begin{quiz}
La funzione $f(x)=\log\left(3+\frac4{x^2}\right)+2x$ ha come asintoto obliquo sinistro la retta
\begin{enumerate}[label=\alph*),leftmargin=1.3cm,itemsep=3pt,topsep=5pt]\item $y=3x+\log3$\item $y=2x+\log4$\item $y=5x+\log4$\item $y=2x+\log3$\item $y=2x$\end{enumerate}
\end{quiz}
% SOLUZIONE: D
% SPIEGAZIONE: Per $x\\to-\\infty$, $4/x^2\\to0$, dunque $\\log(3+4/x^2)=\\log3+o(1)$. Pertanto $f(x)=2x+\\log3+o(1)$ e l'asintoto e $y=2x+\\log3$. Risposta 4 (D).

\begin{quiz}
Sia $f:\mathbb R\to\mathbb R$ una funzione derivabile e strettamente monotona, con $f(4)=5$, $f'(4)=\frac18$. Allora necessariamente
\begin{enumerate}[label=\alph*),leftmargin=1.3cm,itemsep=3pt,topsep=5pt]\item $(f^{-1})'(5)=8$\item $(f^{-1})'(5)=\frac18$\item $(f^{-1})'(5)=4$\item $(f^{-1})'(6)=\frac18$\item $(f^{-1})'(6)=\frac14$\end{enumerate}
\end{quiz}
% SOLUZIONE: A
% SPIEGAZIONE: Derivata dell'inversa: $(f^{-1})'(f(a))=1/f'(a)$ se $f'(a)\\ne0$. Qui $f(4)=5$ e $f'(4)=1/8$, quindi $(f^{-1})'(5)=8$. Risposta 1 (A).

\begin{quiz}
Sia $a_n=(-1)^n\dfrac{\log n}{\sqrt[3]{n}}$. Allora
\begin{enumerate}[label=\alph*),leftmargin=1.3cm,itemsep=3pt,topsep=5pt]\item $a_n$ tende sia a $+\infty$ che a $-\infty$\item esiste $\bar n\in\mathbb N$ tale che $n>\bar n\Rightarrow|a_n|<10^{-8}$\item il limite $\lim_{n\to\infty}a_n$ non esiste\item $a_n$ \`e strettamente crescente\item $a_n$ non \`e limitata superiormente\end{enumerate}
\end{quiz}
% SOLUZIONE: B
% SPIEGAZIONE: Il modulo e $|a_n|=\\log n/n^{1/3}\\to0$: qualsiasi potenza positiva cresce piu velocemente del logaritmo. Anche moltiplicando per $(-1)^n$ il limite rimane $0$. Per definizione di limite, per $\\varepsilon=10^{-8}$ esiste una soglia oltre la quale $|a_n|<\\varepsilon$. Risposta 2 (B).

\par\noindent\hfill{\color{cyan!65!black}\scriptsize Pag. originale 7/521}\par\smallskip
\begin{quiz}
Il limite $\displaystyle\lim_{x\to0}\frac{\sin^2(3x)}{1-\cos(2x)}$ vale
\begin{enumerate}[label=\alph*),leftmargin=1.3cm,itemsep=3pt,topsep=5pt]\item $\frac94$\item $\frac32$\item $\frac92$\item $0$\item $6$\end{enumerate}
\end{quiz}
% SOLUZIONE: C
% SPIEGAZIONE: Usa gli equivalenti notevoli $\\sin(3x)\\sim3x$ e $1-\\cos(2x)\\sim(2x)^2/2=2x^2$. Il rapporto tende a $9x^2/(2x^2)=9/2$. Risposta 3 (C).

\begin{quiz}
Sia $f$ una funzione continua su $[0,2]$ e derivabile in $(0,2)$. Sia inoltre $f(0)<f(2)$. È vero che:
\begin{enumerate}[label=\alph*),leftmargin=1.3cm,itemsep=3pt,topsep=5pt]\item $f$ \`e crescente su $[0,2]$\item $f$ ha un punto di massimo $\bar x\in[0,2]$ tale che $f(\bar x)\ge f(2)$\item $f(2)$ \`e il massimo di $f$ in $[0,2]$\item $f'(x)>0$ per ogni $x\in(0,2)$\item per ogni $x\in[0,2]$, $f(0)\le f(x)\le f(2)$\end{enumerate}
\end{quiz}
% SOLUZIONE: B
% SPIEGAZIONE: Per Weierstrass, la continuita su $[0,2]$ garantisce un massimo assoluto in un punto $\\bar x$. Per definizione $f(\\bar x)\\ge f(2)$. La sola condizione $f(0)<f(2)$ non impone monotonia ne massimo in $2$. Risposta 2 (B).

\begin{quiz}
Sia $f$ una funzione definita su $(-3,3)$ e tale che $\lim_{x\to2}f(x)=-1$. È vero che:
\begin{enumerate}[label=\alph*),leftmargin=1.3cm,itemsep=3pt,topsep=5pt]\item esiste $\delta>0$ tale che $|x-2|<\delta\Rightarrow f(x)<0$\item $f$ \`e continua in $x=2$\item $f(x)<0$\item $f(2)=-1$\item esiste $\delta>0$ tale che $0<|x-2|<\delta\Rightarrow f(x)<0$\end{enumerate}
\end{quiz}
% SOLUZIONE: E
% SPIEGAZIONE: Dalla definizione di limite, poiche $f(x)\\to-1<0$, scegli una tolleranza $\\varepsilon=1/2$: per $0<|x-2|<\\delta$ risulta $-3/2<f(x)<-1/2<0$. Il valore $f(2)$ puo invece essere diverso da $-1$: percio va escluso $x=2$. Risposta 5 (E).

\begin{quiz}
Data $f$ localmente integrabile su $[0,+\infty)$, si dice che l'integrale improprio $\int_0^{+\infty}f(x)\,dx$ \`e divergente se:
\begin{enumerate}[label=\alph*),leftmargin=1.3cm,itemsep=3pt,topsep=5pt]\item $\lim_{t\to+\infty}\int_0^t f(x)\,dx$ \`e uguale a $+\infty$ o a $-\infty$\item $\lim_{t\to+\infty}\int_0^t f(x)\,dx$ non \`e un valore finito\item la successione $a_n=\int_0^n f(x)\,dx$ soddisfa $\lim_{n\to+\infty}|a_n|=+\infty$\item la successione $a_n=\int_0^n f(x)\,dx$ soddisfa $\lim_{n\to+\infty}a_n=+\infty$ o $-\infty$\item $\lim_{t\to+\infty}\left|\int_0^t f(x)\,dx\right|=+\infty$\end{enumerate}
\end{quiz}
% SOLUZIONE: A
% SPIEGAZIONE: Per definizione, un integrale improprio diverge a $+\\infty$ o $-\\infty$ se il limite delle primitive troncate $I(t)=\\int_0^t f(x)\\,dx$ vale rispettivamente $+\\infty$ o $-\\infty$. Se il limite non esiste nemmeno in senso esteso, l'integrale e indeterminato/oscillante, non divergente nel senso richiesto. Risposta 1 (A).

\begin{quiz}
L'insieme delle soluzioni della disequazione $\dfrac{x\log(x+2)}{x-3}\le0$ \`e:
\begin{enumerate}[label=\alph*),leftmargin=1.3cm,itemsep=3pt,topsep=5pt]\item $[-1,0]$\item $(-2,-1]\cup[0,3)$\item nessuna delle altre\item $(3,+\infty)$\item $(-2,+\infty)$\end{enumerate}
\end{quiz}
% SOLUZIONE: B
% SPIEGAZIONE: Dominio: $x>-2$ e $x\\ne3$. Zeri del numeratore: $x=0$ e $\\log(x+2)=0\\iff x=-1$. Segni: su $(-2,-1)$ il prodotto $x\\log(x+2)>0$ e il denominatore e negativo; su $(-1,0)$ il rapporto e positivo; su $(0,3)$ il rapporto e negativo. Includi gli zeri e escludi $3$: $(-2,-1]\\cup[0,3)$. Risposta 2 (B).

\par\noindent\hfill{\color{cyan!65!black}\scriptsize Pag. originale 8/521}\par\smallskip
\begin{quiz}
Sia $f(x)=\sin(\log|x|)$. Questa funzione \`e
\begin{enumerate}[label=\alph*),leftmargin=1.3cm,itemsep=3pt,topsep=5pt]\item pari\item dispari\item non ha alcuna simmetria\item periodica\item crescente\end{enumerate}
\end{quiz}
% SOLUZIONE: A
% SPIEGAZIONE: Verifica la parita: il dominio $\\mathbb R\\setminus\\{0\\}$ e simmetrico e $f(-x)=\\sin(\\log|-x|)=\\sin(\\log|x|)=f(x)$. Risposta 1 (A).

\begin{quiz}
Sia $f:A\subseteq\mathbb R\to\mathbb R$. Si dice punto di minimo per $f$ su $A$:
\begin{enumerate}[label=\alph*),leftmargin=1.3cm,itemsep=3pt,topsep=5pt]\item un punto $x_0\in A$ tale che per ogni $x\in A$, $f(x)>f(x_0)$\item un punto $(x_0,f(x_0))$ tale che per ogni $x\in A$, $f(x)>f(x_0)$\item un punto $(x_0,f(x_0))$ tale che per ogni $x\in A$, $f(x)\ge f(x_0)$\item un punto $f(x_0)$, con $x_0\in A$, tale che per ogni $x\in A$, $f(x)\ge f(x_0)$\item un punto $x_0\in A$ tale che per ogni $x\in A$, $f(x)\ge f(x_0)$\end{enumerate}
\end{quiz}
% SOLUZIONE: E
% SPIEGAZIONE: Definizione di punto di minimo assoluto: $x_0\\in A$ e $f(x_0)\\le f(x)$ per ogni $x\\in A$. Il punto di minimo e l'argomento $x_0$, mentre $f(x_0)$ e il valore minimo e $(x_0,f(x_0))$ e un punto del grafico. Risposta 5 (E).

\begin{quiz}
Sia $f$ una funzione di classe $C^3(\mathbb R)$, tale che $f(x)=-1+2x^2-3x^3+o(x^3)$. È vero che:
\begin{enumerate}[label=\alph*),leftmargin=1.3cm,itemsep=3pt,topsep=5pt]\item $f(1)=-2$\item $f(0)=-1$, $f'(0)=0$\item $f(0)=-1$, $f''(0)=2$\item $\lim_{x\to+\infty}f(x)=-\infty$\item $f'''(0)=3$\end{enumerate}
\end{quiz}
% SOLUZIONE: B
% SPIEGAZIONE: Confronta con Taylor in $0$: $f(x)=f(0)+f'(0)x+f''(0)x^2/2+f'''(0)x^3/6+o(x^3)$. Da $-1+2x^2-3x^3$ segue $f(0)=-1$, $f'(0)=0$, $f''(0)=4$, $f'''(0)=-18$. Risposta 2 (B).

\begin{quiz}
Sia $f$ una funzione definita su un intervallo aperto $I$. Una primitiva di $f$ su $I$ \`e una funzione $F$ tale che:
\begin{enumerate}[label=\alph*),leftmargin=1.3cm,itemsep=3pt,topsep=5pt]\item $F(x)=\int_{x_0}^x f(t)\,dt$, se $x,x_0\in I$\item $F$ \`e definita su $I$ e $F'(x)=f(x)$ per ogni $x\in I$\item $F$ \`e derivabile su $I$ e $F'(x)=f(x)$ per ogni $x\in I$\item $f'(x)=F(x)$ per ogni $x\in I$\item $F$ \`e derivabile in $x_0\in I$ e $F'(x_0)=f(x_0)$\end{enumerate}
\end{quiz}
% SOLUZIONE: C
% SPIEGAZIONE: Per definizione, $F$ e una primitiva di $f$ su $I$ se e derivabile in ogni punto di $I$ e soddisfa $F'(x)=f(x)$ per ogni $x\\in I$. La sola uguaglianza formale senza esplicitare la derivabilita non e la definizione completa. Risposta 3 (C).

\begin{quiz}
La funzione $f(x)=\sqrt[3]{x^2}$
\begin{enumerate}[label=\alph*),leftmargin=1.3cm,itemsep=3pt,topsep=5pt]\item \`e derivabile solo per $x>0$\item ha derivata $f'(x)=\sqrt[3]{2x}$\item ha derivata $f'(x)=\frac23\sqrt{x}$\item non \`e derivabile per $x=0$\item non \`e continua per $x=0$\end{enumerate}
\end{quiz}
% SOLUZIONE: D
% SPIEGAZIONE: Scrivi $f(x)=|x|^{2/3}$. Per $x\\ne0$, $f'(x)=(2/3)\\operatorname{sgn}(x)|x|^{-1/3}$. In $0$ il rapporto incrementale e $|h|^{2/3}/h=\\operatorname{sgn}(h)|h|^{-1/3}$: tende a $+\\infty$ da destra e $-\\infty$ da sinistra. Non e derivabile in $0$, ma e continua. Risposta 4 (D).

\par\noindent\hfill{\color{cyan!65!black}\scriptsize Pag. originale 9/521}\par\smallskip
\begin{quiz}
La funzione $F(x)=\displaystyle\int_x^0 e^{\sin^2 t}\,dt$ \`e tale che:
\begin{enumerate}[label=\alph*),leftmargin=1.3cm,itemsep=3pt,topsep=5pt]\item $F(x)$ \`e crescente, per ogni $x\in\mathbb R$\item $F(x)$ non \`e definita su $\mathbb R$\item $F(x)\le0$, per ogni $x\in\mathbb R$\item $F(x)\ge0$, per ogni $x\in\mathbb R$\item $F(x)$ \`e decrescente, per ogni $x\in\mathbb R$\end{enumerate}
\end{quiz}
% SOLUZIONE: E
% SPIEGAZIONE: Teorema fondamentale del calcolo: $\\frac{d}{dx}\\int_x^0 e^{\\sin^2t}\\,dt=-e^{\\sin^2x}<0$ per ogni $x$. Dunque $F$ e strettamente decrescente su $\\mathbb R$. Risposta 5 (E).

\begin{quiz}
Sia $f(x)=e^x\sqrt{1+e^x}$. Una primitiva di $f(x)$ su $\mathbb R$ \`e:
\begin{enumerate}[label=\alph*),leftmargin=1.3cm,itemsep=3pt,topsep=5pt]\item $F(x)=\frac23\sqrt{(1+e^x)^3}$\item $F(x)=\frac32(1+e^x)^{2/3}$\item $F(x)=\frac1{\sqrt{1+e^x}}$\item $F(x)=e^x\sqrt{1+e^x}$\item $F(x)=\frac{e^x}{\sqrt{1+e^x}}$\end{enumerate}
\end{quiz}
% SOLUZIONE: A
% SPIEGAZIONE: Sostituzione immediata: poni $u=1+e^x$, quindi $du=e^x\\,dx$. Allora $\\int e^x\\sqrt{1+e^x}\\,dx=\\int u^{1/2}du=(2/3)u^{3/2}+C=(2/3)\\sqrt{(1+e^x)^3}+C$. Risposta 1 (A).

\begin{quiz}
La funzione $f(x)=\cos(\sin x)-1+\frac12x^2$ ha sviluppo di Maclaurin all'ordine 2 uguale a:
\begin{enumerate}[label=\alph*),leftmargin=1.3cm,itemsep=3pt,topsep=5pt]\item $o(x^2)$, per $x\to0$\item $x^2+o(x^2)$, per $x\to0$\item non ha sviluppo di Maclaurin di ordine 2\item $0$\item $-2+x^2+o(x^2)$, per $x\to0$\end{enumerate}
\end{quiz}
% SOLUZIONE: A
% SPIEGAZIONE: Sviluppi in $0$: $\\sin x=x+O(x^3)$ e $\\cos t=1-t^2/2+O(t^4)$. Quindi $\\cos(\\sin x)=1-x^2/2+O(x^4)$; sommando $-1+x^2/2$ resta $O(x^4)=o(x^2)$. Risposta 1 (A).

\begin{quiz}
Siano $f(x)\sim\sin x$ e $g(x)\sim x^2$, per $x\to0$. È vero che:
\begin{enumerate}[label=\alph*),leftmargin=1.3cm,itemsep=3pt,topsep=5pt]\item $\lim_{x\to0}\frac{f(x)}{g(x)}=+\infty$\item $f(x)+g(x)=x^2+o(x^2)$, per $x\to0$\item $f(x)+g(x)=x+o(x)$, per $x\to0$\item $f(x)+g(x)=\sin x+x^2$\item non ci sono dati sufficienti per calcolare $\lim_{x\to0}f(x)g(x)$\end{enumerate}
\end{quiz}
% SOLUZIONE: C
% SPIEGAZIONE: Da $f(x)\\sim\\sin x\\sim x$ e $g(x)\\sim x^2$ segue $f(x)=x+o(x)$, $g(x)=x^2+o(x^2)=o(x)$. Sommando: $f(x)+g(x)=x+o(x)$. Risposta 3 (C).

\begin{quiz}
L'enunciato corretto \`e:
\begin{enumerate}[label=\alph*),leftmargin=1.3cm,itemsep=3pt,topsep=5pt]\item se $f$ \`e definita su $[a,b]$ e $f(a)f(b)<0$, allora $f$ ha almeno uno zero in $(a,b)$\item se $f$ \`e continua su $[a,b]$ e $f(a)f(b)<0$, allora $f$ ha almeno uno zero in $(a,b)$\item se $f$ \`e continua su $(a,b)$ e $f(a)f(b)<0$, allora $f$ ha almeno uno zero in $(a,b)$\item se $f$ \`e continua su $[a,b]$ e $f(a)f(b)<0$, allora $f$ ha un unico zero in $(a,b)$\item se $f$ \`e continua su $[a,b]$ e $f(a)f(b)<0$, allora $0\in(a,b)$\end{enumerate}
\end{quiz}
% SOLUZIONE: B
% SPIEGAZIONE: Teorema degli zeri (Bolzano): se $f$ e continua su tutto $[a,b]$ e $f(a)f(b)<0$, allora esiste almeno un $c\\in(a,b)$ con $f(c)=0$. Non si puo affermare che sia unico senza ulteriori ipotesi. Risposta 2 (B).

